\section{Toward a general variational method for random evolution equations}
\label{S:theory-outlook}

The preceding section asks how to extend the results. Here we ask what the
proof itself teaches us. Several long arguments preserve the same three
objects: the displacement--velocity evolution, the covariance of the original
noise, and a variational reward for deterministic amplification. Their
connection suggests a possible method for other random evolution equations.
These questions arose by asking Amy to revisit the calculations and failed
approaches, then checking her proposed abstractions against the arguments.
The present book supplies a detailed case study; a general theorem would
require assumptions that can be verified independently in a second model.

\Needspace{5\baselineskip}
\paragraph{Compose the evolution and retain the joint Gaussian cost.}
For a wave equation, displacement at the end of a time block does not determine
the next block: velocity is also needed. Lemma~\ref{L:raw-block-restart}
composes the full phase-space propagators and identifies them through limits
of the same regularized noise. The Gaussian change of
law in Proposition~\ref{P:raw-block-gaussian-law} acts on that common noise.
It does not replace correlated blocks by independent copies.

Even two Gaussian coordinates show why this distinction matters. If their
covariance matrix is
\[
  C_\rho=\begin{pmatrix}1&\rho\\\rho&1\end{pmatrix},\qquad |\rho|<1,
\]
then a shift by $h=(q,q)$ has Cameron--Martin cost
\[
  \frac12 h^{\mathsf T}C_\rho^{-1}h=\frac{q^2}{1+\rho}.
\]
Indeed, $C_\rho h=(1+\rho)h$ and $|h|^2=2q^2$. The cost is $q^2$ for
independent coordinates and tends to $q^2/2$ as $\rho\uparrow1$.
Discarding correlation can therefore change the leading variational cost.
The Gaussian shift formula is classical~\cite{janson:97:gaussian}*{Theorem~14.1};
the work here is to connect it to the physical wave evolution and its terminal
observable, with estimates uniform in the necessary limits.

This suggests a precise task: find sufficient conditions under which a
phase-space evolution, controlled by one Gaussian field, has matching
long-time upper and lower moment bounds with the same deterministic reward
and the same covariance cost. The conditions must include convergence from
the original regularized equation, control of products on growing intervals,
and passage from phase-space estimates to the chosen observable.
Proposition~\ref{P:clock-copy-obstruction} rules out one unweighted
constant-correlation Brownian-clock substitution; it does not rule out every
probabilistic representation. The pathwise composition used here supplies
neither a conditional-expectation identity nor a general martingale calculus.

\paragraph{Separate the finite optimization from the wave estimates.}
Propositions~\ref{P:gram-positive-optimizer}
and~\ref{P:gram-entrywise-comparison} already apply in any real Hilbert space.
For the square-root reward in~\eqref{E:finite-gram-score}, nonnegative Gram
entries permit a positive conic maximizer and an entrywise comparison of
optimized values. Singular matrices and inactive labels are included. This
is the most explicit piece of the argument that can be stated without wave
notation. Its physical application still needs the covariance and spectral
estimates in Section~\ref{S:intrinsic-second-moment}.

A useful test would replace the square-root reward by the growth rate of a
damped equation. For the scalar equation $u''+2\eta u'=vu$, with $\eta>0$
and $v\ge0$, substituting $u(t)=e^{rt}$ gives
$r^2+2\eta r-v=0$. Its larger root is
\[
  r_\eta(v)=\sqrt{\eta^2+v}-\eta.
\]
Unlike $\sqrt v$, this reward has a finite derivative at zero and is not
homogeneous in $v$ when $\eta$ is fixed. The active-set analysis and dilation
argument therefore need to be checked again. A damped-wave transfer would
also require new bounds for the physical propagator and noise approximation;
the scalar calculation alone supplies none of these bounds.

\paragraph{From an optimized value to the shape of large events.}
The constant $C_2$ identifies an optimized reward, but does not show that the
noise or solution, conditioned on large amplification, is close to an
optimizing profile. Conjecture~\ref{Conj:positive-control-compactness} asks
for compactness of one optimizing sequence after spatial translation.
A further question is whether a suitable conditioned or moment-weighted law
concentrates near such profiles. This needs a specified observable and
scaling, tightness of the relevant laws, and upper bounds away from the
optimizing set. Neither the moment limit nor the compactness conjecture alone
provides these ingredients. This question concerns the shape of rare large
events; its relation to typical fluctuations remains open, as discussed in
Section~\ref{S:future-universality}.

The first test of a broader method should be a second model for which one
abstract statement removes a substantial part of the proof. A damped wave
with correlated Gaussian forcing is a candidate. The Gaussian geometry and
finite optimization should be separated from the estimates specific to its
propagator. If those estimates require rebuilding nearly every argument,
then the current work offers an instructive organization of a proof, while
the proposed general method remains to be found.
